%!TEX root=../../note.tex
% Source: MATH 147 lecture notes; limits and order, squeeze theorem, monotone convergence.
\subsection{Limits, Order, and Monotone Sequences}
Limits preserve order, and completeness turns a bounded monotone sequence into a convergent one. These facts let us find limits even when the limit laws alone do not apply.

\subsubsection{Limits and Inequalities}
\begin{theorem}\label{thm:limit-nonneg}
    If $\lim_{n\to\infty}x_n=a$ and $x_n\geq0$ for every $n\in\N$, then $a\geq0$.\index{limit!preserves order}
    \insymbols{\bigl[\lim_{n\to\infty}x_n=a\wedge(\forall n\in\N)\,[x_n\geq0]\bigr]\Rightarrow a\geq0}
\end{theorem}
\begin{idea}
    If $a<0$, convergence would force every sufficiently late term into a small interval around $a$ that lies below zero.
\end{idea}
\begin{proof}
    Suppose instead that $a<0$. Take $\epsilon=-a/2>0$. Since $\lim_{n\to\infty}x_n=a$, some $n_\epsilon\in\N$ satisfies $\abs{x_n-a}<\epsilon$ whenever $n\geq n_\epsilon$. For such $n$,
    \[
        x_n<a+\epsilon=\frac a2<0,
    \]
    contradicting $x_n\geq0$. Thus $a\geq0$.
\end{proof}

The same conclusion holds if $x_n\geq0$ only for all sufficiently large $n$: the proof uses just one term after both thresholds.

\begin{theorem}\label{thm:limits-and-order}
    If $\lim_{n\to\infty}x_n=a$, $\lim_{n\to\infty}y_n=b$, and $x_n\leq y_n$ for every $n\in\N$, then $a\leq b$.
    \insymbols{\bigl[\lim_{n\to\infty}x_n=a\wedge \lim_{n\to\infty}y_n=b\wedge(\forall n\in\N)\,[x_n\leq y_n]\bigr]\Rightarrow a\leq b}
\end{theorem}
\begin{proof}
    Set $z_n=y_n-x_n\geq0$. The difference law (\cref{thm:limit-laws}, part (iv)) gives $\lim_{n\to\infty}z_n=b-a$. By \cref{thm:limit-nonneg}, $b-a\geq0$, so $a\leq b$.
\end{proof}

Strict inequalities are not preserved: $\frac1n>0$ for every $n\in\N$, yet $\lim_{n\to\infty}\frac1n=0$. From $x_n<y_n$ for all $n$ we may conclude only $a\leq b$.

\subsubsection{The Squeeze Theorem}
\begin{theorem}[Squeeze Theorem]\label{thm:squeeze}
    Suppose $x_n\leq y_n\leq z_n$ for every $n\in\N$ and $\lim_{n\to\infty}x_n=a$, $\lim_{n\to\infty}z_n=a$. Then $\lim_{n\to\infty}y_n=a$.\index{squeeze theorem}
    \insymbols{\bigl[(\forall n\in\N)\,[x_n\leq y_n\leq z_n]\wedge \lim_{n\to\infty}x_n=a\wedge \lim_{n\to\infty}z_n=a\bigr]\Rightarrow \lim_{n\to\infty}y_n=a}
\end{theorem}
\begin{proof}
    Let $\epsilon>0$. Choose $n_1,n_2\in\N$ such that $\abs{x_n-a}<\epsilon$ for $n\geq n_1$ and $\abs{z_n-a}<\epsilon$ for $n\geq n_2$. If $n\geq\max\set{n_1,n_2}$, then
    \[
        a-\epsilon<x_n\leq y_n\leq z_n<a+\epsilon.
    \]
    Thus $\abs{y_n-a}<\epsilon$ for every such $n$, so $\lim_{n\to\infty}y_n=a$.
\end{proof}

\begin{example}\label{ex:sin-over-sqrt}
    Prove that $\lim_{n\to\infty}\frac{\sin n}{\sqrt n}=0$.
\end{example}
\begin{proof}
    Since $-1\leq\sin n\leq1$ and $\sqrt n>0$,
    \[
        -\frac1{\sqrt n}\leq\frac{\sin n}{\sqrt n}\leq\frac1{\sqrt n}.
    \]
    The outer sequences tend to $0$: given $\epsilon>0$, the Archimedean Property (\cref{def:archimedean}) supplies $n_\epsilon\in\N$ with $n_\epsilon>1/\epsilon^2$, and every $n\geq n_\epsilon$ gives $1/\sqrt n\leq1/\sqrt{n_\epsilon}<\epsilon$; the lower bound $-1/\sqrt n$ tends to $0$ by the same $n_\epsilon$. \Cref{thm:squeeze} now gives the result.
\end{proof}

\subsubsection{Monotone Sequences}
\begin{definition}\label{def:monotone}
    A sequence $\set{x_n}$ is \emph{increasing} if $x_n\leq x_{n+1}$ for every $n\in\N$, and \emph{decreasing} if $x_n\geq x_{n+1}$ for every $n\in\N$. It is \emph{monotone} if it is increasing or decreasing.\index{monotone sequence}
\end{definition}

% Keep the statement and its symbolic form in one frame.
\mdfsetup{nobreak=true}
\begin{theorem}[Monotone Convergence Theorem]\label{thm:monotone-convergence}
    A monotone sequence of real numbers converges if and only if it is bounded. More precisely:\index{monotone convergence theorem}
    \begin{enumerate}
        \item[(i)] If $\set{x_n}$ is bounded and increasing, then
            \[
                \lim_{n\to\infty}x_n=\sup\set{x_n:n\in\N}.
            \]
        \item[(ii)] If $\set{x_n}$ is bounded and decreasing, then
            \[
                \lim_{n\to\infty}x_n=\inf\set{x_n:n\in\N}.
            \]
    \end{enumerate}
    \insymbols{\begin{aligned}
        &[\set{x_n}\text{ monotone}]\\
        &\quad\Rightarrow
            \bigl[(\exists a\in\R)\,[\lim_{n\to\infty}x_n=a]\iff\set{x_n}\text{ bounded}\bigr],\\
        \text{(i)}\quad &[\set{x_n}\text{ bounded and increasing}]\\
        &\quad\Rightarrow\lim_{n\to\infty}x_n=\sup\set{x_n:n\in\N},\\
        \text{(ii)}\quad &[\set{x_n}\text{ bounded and decreasing}]\\
        &\quad\Rightarrow\lim_{n\to\infty}x_n=\inf\set{x_n:n\in\N}.
    \end{aligned}}
\end{theorem}
\mdfsetup{nobreak=false}
\begin{idea}
    Completeness gives a supremum or infimum. A term lies within $\epsilon$ of that bound, and monotonicity keeps every later term there.
\end{idea}
\begin{proof}
    Every convergent sequence is bounded (\cref{thm:convergent-bounded}), proving the forward implication. For the converse, (i) and (ii) cover the two kinds of monotone sequence.

    For (i), suppose $\set{x_n}$ is bounded and increasing. The set $S=\set{x_n:n\in\N}$ is nonempty and bounded above, so completeness gives $s=\sup S$. Let $\epsilon>0$. The $\epsilon$-characterization of the supremum (\cref{lem:eps-characterization}) gives an index $n_\epsilon$ with $s-\epsilon<x_{n_\epsilon}$. For every $n\geq n_\epsilon$, monotonicity and the upper-bound property give
    \[
        s-\epsilon<x_{n_\epsilon}\leq x_n\leq s<s+\epsilon.
    \]
    Hence $\lim_{n\to\infty}x_n=s$.

    For (ii), suppose $\set{x_n}$ is bounded and decreasing. The set $S$ is nonempty and bounded below, so $t=\inf S$ exists (\cref{cor:inf-existence}). Let $\epsilon>0$. Since $t+\epsilon>t$ and $t$ is the \emph{greatest} lower bound, $t+\epsilon$ is not a lower bound of $S$; so there is an index $n_\epsilon$ with $x_{n_\epsilon}<t+\epsilon$. For $n\geq n_\epsilon$, monotonicity and the lower-bound property give
    \[
        t-\epsilon<t\leq x_n\leq x_{n_\epsilon}<t+\epsilon.
    \]
    Thus $\lim_{n\to\infty}x_n=t$.
\end{proof}

\subsubsection{A Recursively Defined Sequence}
The Monotone Convergence Theorem is most useful for sequences defined by recursion, where no formula for $x_n$ is available. It shows that a limit exists; the recursion then determines its value.

\begin{example}\label{ex:recursive-mct}
    Let $x_1=1$ and $x_{n+1}=\frac{2x_n+3}{4}$ for all $n\in\N$. Show that $\set{x_n}$ converges and find its limit.
\end{example}
The first terms are
\[
    x_1=1,\qquad x_2=\frac{2\cdot1+3}{4}=\frac54=1.25,\qquad x_3=\frac{2\cdot\frac54+3}{4}=\frac{11}8=1.375,
\]
which suggests that the sequence increases and stays below $2$.
\begin{idea}
    Prove by induction that the sequence is bounded above and increasing; the Monotone Convergence Theorem then gives a limit $a$, and letting $n\to\infty$ in the recursion gives an equation for $a$.
\end{idea}
\begin{proof}
    \textbf{Step 1: $x_n<2$ for all $n\in\N$.} We use induction on $n$. For $n=1$, $x_1=1<2$. Let $k\in\N$ and assume $x_k<2$. Then
    \[
        2x_k<4\implies 2x_k+3<7\implies x_{k+1}=\frac{2x_k+3}{4}<\frac74<2.
    \]
    By induction, $x_n<2$ for every $n\in\N$.

    \textbf{Step 2: $x_n\leq x_{n+1}$ for all $n\in\N$.} Again by induction. For $n=1$, $x_1=1<\frac54=x_2$. Let $k\in\N$ and assume $x_k\leq x_{k+1}$. Then
    \[
        2x_k+3\leq 2x_{k+1}+3\implies x_{k+1}=\frac{2x_k+3}{4}\leq\frac{2x_{k+1}+3}{4}=x_{k+2}.
    \]
    By induction, $\set{x_n}$ is increasing.

    \textbf{Step 3: the limit.} By Step 2, $x_n\geq x_1=1$ for every $n$, and by Step 1, $x_n<2$; so $\set{x_n}$ is bounded and increasing. By \cref{thm:monotone-convergence} there is $a\in\R$ with $\lim_{n\to\infty}x_n=a$.

    The shifted sequence also satisfies $\lim_{n\to\infty}x_{n+1}=a$: if $\abs{x_n-a}<\epsilon$ for all $n\geq n_\epsilon$, then also $\abs{x_{n+1}-a}<\epsilon$ for all $n\geq n_\epsilon$, since $n+1\geq n_\epsilon$. On the other hand, $x_{n+1}=\frac12x_n+\frac34$, so the limit laws (\cref{thm:limit-laws}, parts (i)--(iii)) give
    \[
        \lim_{n\to\infty}x_{n+1}=\frac12a+\frac34=\frac{2a+3}{4}.
    \]
    A limit is unique (\cref{thm:limit-unique}), so
    \[
        a=\frac{2a+3}{4}\implies 4a=2a+3\implies a=\frac32.
    \]
    Hence $\lim_{n\to\infty}x_n=\frac32$.
\end{proof}

Existence must come first. The equation $a=\frac{2a+3}{4}$ only says what the limit would be \emph{if} there is one: for $x_1=1$ and $x_{n+1}=2x_n$, the same manipulation gives $a=2a$, so $a=0$, yet $x_n=2^{n-1}$ diverges to $+\infty$. The bound $2$ in Step 1 is not special; any upper bound at least $\frac32$ (the solution of the limit equation) would work.

\subsubsection*{In practice}
\begin{itemize}
    \item \textbf{Use the squeeze theorem.} Bound a difficult sequence above and below by sequences with the same known limit; verify that both bounds hold for all sufficiently large indices.
    \item \textbf{Use monotonicity.} Prove the direction of monotonicity and find a bound. Completeness then gives the limit as the supremum or infimum of the terms.
    \item \textbf{Recursive sequences.} Compute a few terms to guess the direction and a bound; prove both by induction; apply the Monotone Convergence Theorem; then pass to the limit on both sides of the recursion and solve for the limit.
\end{itemize}
